\documentclass[11pt]{article}
\usepackage{fontspec}
\setmainfont{Times New Roman}
\setsansfont{Arial}
\setmonofont{Consolas}
\usepackage{amsmath,amssymb}
\usepackage{geometry}
\usepackage{microtype}
\geometry{a4paper,margin=3cm}
\setlength{\parindent}{1.25em}
\setlength{\parskip}{0pt}
\emergencystretch=14em
\allowdisplaybreaks
\providecommand{\vdotswithin}[1]{\mathrel{\vdots}}

\begin{document}

\section*{\S\ 1. Dokaz bazovyh svojstv 1) i 2) za vse oblasti $\mathfrak G$.}

Nehaj $\mathfrak G$ bude koja-koli zadana oblast iz cělyh transcendentnyh čisel, ktorej tedy naležet abstraktna svojstva I--IV. Po postupku, priměnjenom od E. Zermela na polje vsih kompleksnyh čisel, vybirajemo iz $\mathfrak G$ {\itshape algebraičnu bazu} $H$ {\itshape vsih čisel iz $\mathfrak G$}. Poněže po II {\itshape vsako} realno ili kompleksno čislo možno predstaviti kako kvocient dvoh čisel iz $\mathfrak G$, $H$ jest zároveň algebraična baza {\itshape vsih} čisel. Zato za vsaku zadanu oblast $\mathfrak G$ možno izbrati algebraičnu bazu vsih čisel tako, že ona naleži k $\mathfrak G$; tym samym {\itshape bazovo svojstvo 1) jest izpolnjeno}. Dalje, po III, $\mathfrak G$ soderži oblast cělosti $[K]$ vsih cělyh algebraičnyh čisel, i zato po I takože vse polinomy od bazovyh veličin $\eta$ s koeficientami iz $[K]$, t. j. vse celočislene polinomy od $\eta$, ktore obrazuju oblast cělosti $[H]$. Tako {\itshape bazovo svojstvo 2) jest takože vsegda izpolnjeno}; t. j. {\itshape vsaka oblast $\mathfrak G$ može biti skonstruovana posrědstvom algebraične bazy $H$ vsih čisel tako, že $\mathfrak G$ soderži oblast cělosti $[H]$ kako dělitelj.}

\textbf{Primětka.} Naravno možno vse-taki, izhodeči od bazy $H$, skonstruovati oblast $\mathfrak G$, ktorej bazove svojstva 1) i 2) naležet ne upravo vzhodno k $H$, ale vzhodno k nekojej drugoj bazě $H'$. Nehaj, napr., $\mathfrak G'$ vznikaje iz Zermelovej oblasti $\mathfrak G_\eta$ tym, že v cěloj konstrukciji někoju bazovu veličinu $\eta$ zaměnjamo polinomom $f(\eta)$. $\mathfrak G'$ ne soderži ni $\eta$, ni $[H]$; ale ako v bazě $H$ zaměnimo bazovu veličinu $\eta$ čislom $\xi=f(\eta)$ --- pri čem $H$ znovu prěhodi v algebraičnu bazu $H'$ ---, togda bazove svojstva 1) i 2) sut izpolnjene za $\mathfrak G'$ vzhodno k $H'$.

\end{document}
